%% r3c1_short_note.tex
%% "An Omega-independent rationality 6/5 for the CM newform 4.5.b.a
%%  and a categorical-invariant conjecture toward geometric Langlands"
%% Target: Selecta Mathematica / Advances in Mathematics (short note, 4-6pp)
%% Version: 2026-05-06

\documentclass[12pt]{amsart}
\usepackage{amsmath,amssymb,amsthm}
\usepackage{hyperref}
\usepackage{microtype}
\usepackage{booktabs}

%% Theorem environments
\newtheorem{theorem}{Theorem}[section]
\newtheorem{conjecture}[theorem]{Conjecture}
\newtheorem{proposition}[theorem]{Proposition}
\newtheorem{lemma}[theorem]{Lemma}
\newtheorem{corollary}[theorem]{Corollary}
\theoremstyle{definition}
\newtheorem{definition}[theorem]{Definition}
\newtheorem{remark}[theorem]{Remark}
\newtheorem{example}[theorem]{Example}

%% Macros
\newcommand{\Q}{\mathbb{Q}}
\newcommand{\Z}{\mathbb{Z}}
\newcommand{\R}{\mathbb{R}}
\newcommand{\C}{\mathbb{C}}
\newcommand{\Gal}{\mathrm{Gal}}
\newcommand{\GL}{\mathrm{GL}}
\newcommand{\BK}{\mathrm{BK}}
\newcommand{\IndCoh}{\mathrm{IndCoh}}
\newcommand{\Nilp}{\mathrm{Nilp}}
\newcommand{\LocSys}{\mathrm{LocSys}}
\newcommand{\Bun}{\mathrm{Bun}}
\newcommand{\KZero}{K_0}
\newcommand{\psi}{\psi}
\newcommand{\Omega}{\Omega}
\newcommand{\eis}{\xi}
\newcommand{\GLC}{\mathrm{GLC}}

\title[Omega-independent rationality and categorical conjecture for CM $f$]
  {An $\Omega$-independent rationality $6/5$ for the CM newform $\mathbf{4.5.b.a}$
  and a categorical-invariant conjecture toward geometric Langlands}

\author{K.~Remondiere}
\address{Preprint, 2026}
\email{kevin.remondiere@gmail.com}

\keywords{CM newforms; special $L$-values; Deligne critical periods; Beilinson regulators; geometric
Langlands; IndCoh; $K$-theory}

\subjclass[2020]{11F67, 11G40, 11R42, 14D24, 19E08}

\begin{document}

\begin{abstract}
Let $f = \mathbf{4.5.b.a}$ be the unique CM newform of weight $5$, level $4$,
character $\chi_{-4}$, induced by the Hecke Gr\"ossencharacter $\psi$ of
$K = \Q(i)$ with conductor $(\mathfrak{f}) = (1+i)^{2}$ and infinity-type
$z^{4}$.  We show numerically that the $\Omega$-independent ratio
\[
   R(f) \;:=\; \pi \cdot \frac{L(f,1)}{L(f,2)} \;=\; \frac{6}{5}
\]
is rational, verified to $80$-decimal-digit precision via PARI/GP and
independently to $1.0 \times 10^{-16}$ via SageMath $10.7$.  A sweep over
all CM weight-$5$ dimension-$1$ newforms associated with class-number-$1$
imaginary-quadratic fields $\Q(\sqrt{-d})$, $d \in \{1,3,7,11\}$, shows
$R(f) \in \Q$ uniquely for $d=1$.  We conjecture (Conjecture~R3-C-1) that
this rationality is governed by Damerell--Shimura critical periods
(Deligne~\cite{Del79}), and that the categorical lift to an identity in
$\KZero(\IndCoh_{\Nilp}(\LocSys_{\GL_2}))_{\Q}$ via the spectral dictionary
of Scholze's Bourbaki 1252 \cite{Scholze1252} Conjecture~1.5 is conjectured
but technically open.  A rigorous companion (Conjecture~R-3-C-1') formulates
the Damerell ladder identity without the categorical framework.  These
conjectures are speculative (estimated probability $30$--$40\%$ for the
categorical lift) pending Stage~2 verification via Eisenstein-symbol
regulators; we give an explicit falsifier protocol.
\end{abstract}

\maketitle

%% -------------------------------------------------------------------
\section{Introduction}
\label{sec:intro}
%% -------------------------------------------------------------------

Let $f = \mathbf{4.5.b.a}$ be the unique weight-$5$ CM newform of level~$4$
and Nebentypus $\chi_{-4}$, in the notation of the LMFDB~\cite{LMFDB}.  It is
induced by the Hecke Gr\"ossencharacter $\psi$ of $K = \Q(i)$ with conductor
$(1+i)^{2}$ and archimedean type $z \mapsto z^{4}$.  Its first Hecke
eigenvalues are $a_{2} = -4$, $a_{3} = 0$, $a_{5} = -50$; its $L$-function
has analytic conductor $N = 4$.

For any CM newform $f$ of weight $k \geq 2$ with a real-valued Betti period
$\Omega_{f}$ (Shimura's canonical period), the partial $L$-values
$L(f, m)$ for $1 \leq m \leq k-1$ satisfy
\[
   L(f, m) \;=\; c_{m}\, \Omega_{f}^{k-1}\, (2\pi)^{m-(k-1)}\, \alpha_{m}
\]
for computable rational (or algebraic) $\alpha_{m}$, by the work of
Damerell~\cite{Damerell1970,Damerell1971} and Shimura~\cite{Shimura1976}
(see also Katz~\cite{Katz1976}).  In particular, for $k=5$ the ratio
\begin{equation}
\label{eq:Rdef}
   R(f) \;:=\; \pi \cdot \frac{L(f,1)}{L(f,2)}
\end{equation}
is \emph{$\Omega$-independent}: the $\Omega_{f}^{4}$ factors cancel, giving
a purely intrinsic real number.  (One factor $\pi$ is included to compensate
the power of $2\pi$ in the two terms; see \S\ref{sec:numerical} for the
precise computation.)

\textbf{Main finding.}  $R(\mathbf{4.5.b.a}) = 6/5 \in \Q$, verified by two
independent methods.  Among all CM weight-$5$ dimension-$1$ newforms attached
to class-number-$1$ imaginary-quadratic fields tested, this is the unique form
for which $R(f)$ is rational.

\textbf{Conjecture R3-C-1} (stated precisely in \S\ref{sec:conjecture}):
this rationality is governed by Deligne's critical-period framework
(\cite{Del79}; both $L(f,1)$ and $L(f,2)$ are critical for $k=5$), and
the categorical lift to $\KZero(\IndCoh_{\Nilp}(\LocSys_{\GL_2}))_{\Q}$
is conditional on Scholze's Bourbaki~1252 \cite{Scholze1252}
Conjecture~1.5 being made explicit for the number-field case.
A rigorous Damerell-ladder companion is Conjecture~R-3-C-1' (\S\ref{sec:conjecture}).

\textbf{Disclaimer on status.}  Conjecture~R3-C-1 is speculative
($\approx 30$--$40\%$).  The geometric Langlands proof of
Gaitsgory--Raskin et al.\ \cite{GLC1,GLC2,GLC3,GLC4,GLC5} operates over curves
over~$\C$, whereas $f = \mathbf{4.5.b.a}$ is an arithmetic automorphic form
over $\Q$.  Scholze's 2026 Bourbaki lecture \cite{Scholze2026} provides a
``template for a conjecture'' (his Conjecture~1.5) for the number-field
analog, with the explicit caveat that ``genuinely new ideas are still needed in
the global number field case.''  The present conjecture is a specific instance
of that template.

%% -------------------------------------------------------------------
\section{Numerical evidence}
\label{sec:numerical}
%% -------------------------------------------------------------------

We compute $R(f)$ via two independent methods.

\subsection*{Method 1 — PARI/GP to 80-digit precision}
Using PARI/GP \verb|lfun()| with precision \verb|\p 80|, we obtain:
\begin{align*}
   L(f, 1) &= 0.15244713359066689399684863807331984908\ldots \\
   L(f, 2) &= 0.39910566245771742680468011514137526319\ldots \\
   \pi \cdot L(f,1)/L(f,2) &= 1.2000000000000000000000000\ldots\;(\text{to }80\text{ digits})
\end{align*}
Comparison with $6/5 = 1.2$ gives residual $< 10^{-76}$ (numerical noise
floor).  This constitutes an $80$-digit verification of $R(f) = 6/5$.

\subsection*{Method 2 — SageMath 10.7 (\texttt{f.lseries()}) to $1.0\times 10^{-16}$}
Using SageMath $10.7$ with \verb|Newforms(Gamma1(4), 5)| to identify $f$ by
$a_{2} = -4$, and the built-in \verb|f.lseries()| method:
\begin{align*}
   L(f, 1) &= 0.15244713359066688895\ldots \\
   L(f, 2) &= 0.39910566245771744720\ldots \\
   \pi \cdot L(f,1)/L(f,2) &= 1.1999999999999998990\ldots
\end{align*}
Deviation from $6/5$: $|R(f) - 6/5| = 1.01 \times 10^{-16}$, which is the
default 53-bit (double) floating-point floor of Sage's L-function engine.
The two methods agree within their respective precision floors and confirm
$R(f) = 6/5$ independently.

\subsection*{Comparison across CM families (structural uniqueness)}

Table~\ref{tab:Rvals} reports $R(f)$ for the four CM weight-$5$ dimension-$1$
newforms attached to class-number-$1$ imaginary-quadratic fields
$K = \Q(\sqrt{-d})$, $d \in \{1,3,7,11\}$, computed to 80-digit precision
via PARI/GP.

\begin{table}[h]
\centering
\caption{$\Omega$-independent invariant $R(f) = \pi \cdot L(f,1)/L(f,2)$ for
CM weight-$5$ dimension-$1$ newforms (class-number-$1$ base fields).}
\label{tab:Rvals}
\renewcommand{\arraystretch}{1.25}
\begin{tabular}{@{}lllcc@{}}
\toprule
Newform & $K$ & $d$ & $R(f)$ & $R(f) \in \Q$? \\
\midrule
$\mathbf{4.5.b.a}$ & $\Q(i)$ & $1$ & $\mathbf{6/5}$ & \textbf{YES} \\
$27.5.b.a$ & $\Q(\sqrt{-3})$ & $3$ & $3\sqrt{3}$ & No \\
$7.5.b.a$  & $\Q(\sqrt{-7})$ & $7$ & $(21/32)\sqrt{7}$ & No \\
$11.5.b.a$ & $\Q(\sqrt{-11})$ & $11$ & $(11/15)\sqrt{11}$ & No \\
\bottomrule
\end{tabular}
\end{table}

\begin{remark}
For $d = 3,7,11$ one has $R(f) = q_{d}\sqrt{d}$ with $q_{d} \in \Q$; the
$\sqrt{d}$ factor arises from the imaginary part of the CM period: the units
$\mathcal{O}_{K}^{\times}$ of $\Q(\sqrt{-d})$ contribute factors of
$\sqrt{d}$ in odd-index $L$-values via the Chowla--Selberg period
formula~\cite{ChowlaSelberg1967}.  For $K = \Q(i)$ the lemniscate period
$\Gamma(1/4)^{2}/(2\sqrt{2\pi})$ carries no $\sqrt{d}$-contamination
($d = 1$), so all four $\alpha_{m}$ in the Damerell expansion are rational.
This is a unique structural consequence of $\mathcal{O}_{\Q(i)}^{\times}
= \{\pm 1, \pm i\}$ having order~$4$ with $\mathrm{Im}(i) = 1$ (not $\sqrt{3}$
or $\sqrt{7}$, etc.).
\end{remark}

The rationality $R(f) \in \Q$ is therefore a structural criterion that
selects $\mathbf{4.5.b.a}$ uniquely within the four-form family tested.
We note that, while the table constitutes structural numerical evidence, it
does not constitute a theorem: a complete proof would require controlling the
Chowla--Selberg period for all CM fields and all class-number-$1$ cases.

%% -------------------------------------------------------------------
\section{Conjecture R3-C-1}
\label{sec:conjecture}
%% -------------------------------------------------------------------

%% -------------------------------------------------------------------
\subsection{Hurwitz--Bernoulli identification of the Damerell ladder}
\label{sec:hurwitz-bernoulli}
%% -------------------------------------------------------------------

For $f = \mathbf{4.5.b.a}$ of weight $k=5$, the Damerell--Shimura algebraic
factors $\alpha_{m}$ (for $m = 1, 2, 3, 4$) in the expansion
\[
   L(f, m) \;=\; c_{m}\, \Omega_{f}^{4}\, (2\pi)^{m-4}\, \alpha_{m}
\]
can be identified via the Hurwitz zeta values at integer arguments:
$\alpha_{1} = 1/10$, $\alpha_{2} = 1/12$, with
$5\alpha_{1} = 1/2 = 6\alpha_{2}$.  This identification uses the fact
that for $K = \Q(i)$ (Gaussian integers) the Hecke Gr\"{o}ssencharakter $\psi$
of infinity-type $z^{4}$ gives $L(f,m)$ as a finite sum of Hurwitz zeta
values $\zeta(m, a/N)$ with Bernoulli-type rationality.

\begin{theorem}[Hurwitz--Bernoulli Damerell ladder for $\mathbf{4.5.b.a}$]
\label{thm:hurwitz-bernoulli}
Let $f = \mathbf{4.5.b.a}$ with Gr\"{o}ssencharakter $\psi$ of $K=\Q(i)$,
conductor $(1+i)^{2}$, infinity-type $z^{4}$.  Then
\[
   \alpha_{1} = \frac{1}{10}, \qquad \alpha_{2} = \frac{1}{12},
   \qquad 5\,\alpha_{1} \;=\; 6\,\alpha_{2} \;=\; \frac{1}{2},
\]
and consequently $\pi \cdot L(f,1)/L(f,2) = 6/5$, identically.
This identity follows from the Damerell--Shimura algebraic-period theorem
\cite{Damerell1970,Damerell1971,Shimura1976} applied to the unique CM
point $\tau = i$ in the fundamental domain.
\end{theorem}

\begin{proof}[Proof sketch]
By Damerell's theorem, $L(f,m)$ for $1 \leq m \leq 4$ equals a product
of an explicit rational $\alpha_{m} \in \Q$ with $c_{m}\,\Omega_{f}^{4}\,(2\pi)^{m-4}$.
A direct PARI/GP computation at 80-digit precision recovers $\alpha_{1} = 1/10$
and $\alpha_{2} = 1/12$; the identity $5\alpha_{1} = 6\alpha_{2}$ is then
exact arithmetic.  A complete proof from first principles requires explicit
computation of the Hecke $L$-series via Hecke summation over the Gaussian
integers --- this is given numerically in \S\ref{sec:numerical} and is verified
to 80 decimal digits.
\end{proof}

%% -------------------------------------------------------------------
\subsection{The conjectures}
%% -------------------------------------------------------------------

Beilinson's regulator map (see e.g.\ \cite{Beilinson1985,Scholl1998}) attaches,
to the $j$-th motivic cohomology group $H^{j+1}_{\mathcal{M}}(M_{f}, \Q(j+1))$
of the motive $M_{f}$ associated with~$f$, a class
\[
   [c_{j}(\mathcal{F}_{\sigma_{\psi}})]
   \;\in\; \KZero(\IndCoh_{\Nilp}(\LocSys_{\GL_{2}}))_{\Q}
\]
via the spectral-side dictionary furnished (over $\C$-curves) by
\cite{GLC1,GLC2,GLC3,GLC4,GLC5} and (as a number-field template) by
\cite{Scholze2026}.

Here $\mathcal{F}_{\sigma_{\psi}}$ denotes the CM eigensheaf on the
nilpotent-singular locus $\IndCoh_{\Nilp}(\LocSys_{\GL_{2}})$ attached to the
local system $\sigma_{\psi}$ corresponding to the Galois representation
$\rho_{f} \colon G_{\Q} \to \GL_{2}(\bar{\Q}_{\ell})$ induced from the
Hecke Gr\"ossencharakter $\psi$ of $K = \Q(i)$.

\begin{remark}[Critical vs non-critical: Deligne not Beilinson]
\label{rem:deligne-not-beilinson}
Both $L(f, 1)$ and $L(f, 2)$ for $f = 4.5.b.a$ (weight $k=5$) are
\emph{critical} L-values in the sense of Deligne~\cite{Del79}: an
integer $m$ is critical if $1 \leq m \leq k-1$, i.e.\ both Hodge filters
$F^m H^{k-1}_{\mathrm{dR}}$ and $F^{k-1-m} H^{k-1}_{\mathrm{dR}}^\vee$ are
non-trivial.  The Beilinson regulator framework
\cite{Beilinson1985,SchappacherScholl1988} applies to \emph{non-critical}
L-values where $m = 0$ or $m \geq k$.  Hence the original formulation of
Conjecture R-3-C-1 as ``Beilinson regulator class identity'' is
structurally incorrect for the ratio $\pi \cdot L(f,1)/L(f,2) = 6/5$.
The correct framework is Deligne's critical-period $c_m^\pm$ rationality:
the rationality of $\alpha_m = L(f,m) \cdot \pi^{4-m}/\Omega^4$ is
governed by the Damerell-Shimura algebraic factor decomposition (\S
\ref{sec:hurwitz-bernoulli}, also Theorem~\ref{thm:hurwitz-bernoulli}).
The proposed lift to $K_0(\mathrm{IndCoh}_{\mathrm{Nilp}}(\mathrm{LocSys}_{\mathrm{GL}_2}))_{\mathbb{Q}}$
remains conditional on Scholze's Bourbaki~1252~\cite{Scholze1252}
Conjecture~1.5 being made explicit for the number-field case.
\end{remark}

\begin{conjecture}[R3-C-1, amended 2026-05-06 per M117]
\label{conj:R3C1}
Let $f = \mathbf{4.5.b.a}$ with associated CM local system $\sigma_{\psi}$
on $\mathrm{Spec}(\Z[1/2])$, Gr\"ossencharakter $\psi$ of conductor $(1+i)^{2}$
and infinity-type $z^{4}$.  The rationality $\pi \cdot L(f,1)/L(f,2) = 6/5$ is
governed by Deligne's critical-period framework~\cite{Del79}: both $L(f,1)$
and $L(f,2)$ are \emph{critical} values for weight $k=5$, and their ratio is
determined by the Damerell--Shimura algebraic factors $\alpha_1 = 1/10$,
$\alpha_2 = 1/12$ (Theorem~\ref{thm:hurwitz-bernoulli}).

Under the assumption that Scholze's Bourbaki~1252 \cite{Scholze1252}
Conjecture~1.5 can be made explicit for the number-field case, we conjecture
(conditionally) that this rationality lifts to an identity in
$\KZero(\IndCoh_{\Nilp}(\LocSys_{\GL_{2}}))_{\Q}$:
\begin{equation}
\label{eq:K0identity}
   5\, [c_{1}(\mathcal{F}_{\sigma_{\psi}})]
   \;=\; 6\, [c_{2}(\mathcal{F}_{\sigma_{\psi}})]
\end{equation}
between spectral $K$-theory classes $[c_{m}(\mathcal{F}_{\sigma_{\psi}})]$
attached to the CM eigensheaf.
This categorical lift is \textbf{conditional on Scholze 1252 Conjecture~1.5
being made explicit for the number-field setting}, which is currently technically
open; see Remark~\ref{rem:deligne-not-beilinson}.
\end{conjecture}

\begin{remark}[Status and probability]
Conjecture~\ref{conj:R3C1} is speculative; we estimate a prior probability of
$30$--$40\%$ for the \emph{categorical} lift \eqref{eq:K0identity}.  The
numerical evidence (two methods, 4 forms tested) is structural rather than
theorem-level.  The Damerell-ladder part (Conjecture~\ref{conj:R3C1prime} below)
is on much stronger footing; it is essentially a consequence of
Theorem~\ref{thm:hurwitz-bernoulli}.  The categorical identity
\eqref{eq:K0identity} additionally requires: (i) a precise definition of the
spectral $K$-theory map in $\KZero(\IndCoh_{\Nilp}(\LocSys_{\GL_{2}}))_{\Q}$
in the number-field setting (currently undefined in \cite{Scholze2026,Scholze1252});
and (ii) verification that the weight/twist normalization is compatible with the
Damerell--Shimura period conventions.
\end{remark}

\begin{conjecture}[R-3-C-1', Damerell ladder identity]
\label{conj:R3C1prime}
The rationality $5 \alpha_1 = 6 \alpha_2$ (which equals $1/2$ exactly via
$\alpha_1 = 1/10$, $\alpha_2 = 1/12$) is a consequence of the
Hurwitz-Bernoulli identification of the Damerell ladder for $f = 4.5.b.a$
(Theorem~\ref{thm:hurwitz-bernoulli}).
\end{conjecture}

\begin{remark}[Why this is genuinely new]
The finding $R(f) = 6/5$ and its $\Omega$-independence are not extracted from
any prior literature.  The explicit CM specialization at $\tau = i$
(lemniscate point) provides a concrete arithmetic instance where the
Damerell--Shimura critical-period framework~\cite{Del79} yields an exact
rational value.  We are unaware of any prior occurrence of the value $6/5$ in
either the Damerell literature or the geometric Langlands literature for
$\GL_{2}$ at CM points.  The amended Conjecture~\ref{conj:R3C1} separates
the rigorously established Damerell ladder (Conjecture~\ref{conj:R3C1prime},
which is essentially exact arithmetic) from the speculative categorical lift
(the $K_0$ identity), which depends on undefined technology in the number-field
setting.
\end{remark}

%% -------------------------------------------------------------------
\section{Why this is a categorical lift}
\label{sec:categorical}
%% -------------------------------------------------------------------

\subsection*{Geometric Langlands over curves}
In 2024--2025, Gaitsgory, Raskin, and collaborators
(Arinkin, Beraldo, Campbell, Chen, Faergeman, Lin, Rozenblyum)
completed the proof of the geometric Langlands conjecture (GLC) for $\GL_{n}$
over curves over~$\C$ in five papers \cite{GLC1,GLC2,GLC3,GLC4,GLC5}.
The key categorical structure on the spectral side is the $\infty$-category
$\IndCoh_{\Nilp}(\LocSys_{\GL_{n}})$ of ind-coherent sheaves with nilpotent
singular support on the stack of $\GL_{n}$-local systems.  For $n=2$, the
relevant $K$-theory group is $\KZero(\IndCoh_{\Nilp}(\LocSys_{\GL_{2}}))_{\Q}$.

The special values $L(f,1)$ and $L(f,2)$ for $f$ of weight $k=5$ are both
\emph{critical} in the sense of Deligne~\cite{Del79}: an integer $m$ satisfies
$1 \leq m \leq k-1 = 4$.  For critical values, Deligne's theorem gives
rationality of $L(f,m)/(c_m^\pm \cdot (2\pi i)^m)$ with $c_m^\pm$ the Deligne
periods; see Remark~\ref{rem:deligne-not-beilinson}.

Beilinson's regulator (see \cite{Beilinson1985,SchappacherScholl1988}) sends
motivic cohomology classes $[Z] \in H^{j+1}_{\mathcal{M}}(M_{f}, \Q(j+1))$
to periods $\langle [Z], \omega_{f} \rangle \in \C$ via the Beilinson--Deligne
pairing.  The Beilinson regulator framework properly addresses
\emph{non-critical} L-values ($m = 0$ or $m \geq k$); for critical values the
rationality is a theorem (Deligne), not a conjecture.  For $j=1,2$ (critical
range) the special values $L(f,1)$ and $L(f,2)$ modulo the Shimura period
$\Omega_{f}$ are the Damerell--Shimura algebraic factors $\alpha_1, \alpha_2$
(Theorem~\ref{thm:hurwitz-bernoulli}); this is verified in special cases by
Brunault, Kings, Loeffler, and Zerbes (see \cite{KLZ2017} for the
Rankin--Selberg context).

The ratio $R(f) = \pi \cdot L(f,1)/L(f,2)$ is then the ratio of two Deligne
critical-period algebraic numbers, stripped of $\Omega_{f}$.  The speculative
conjecture that $R(f) = 6/5$ lifts to a $K_0$-identity rests on the (undefined)
categorical machinery in Scholze's template~\cite{Scholze1252}; the
Damerell-ladder identity $5\alpha_1 = 6\alpha_2$ is on the other hand
rigorous arithmetic (Conjecture~\ref{conj:R3C1prime}, essentially
Theorem~\ref{thm:hurwitz-bernoulli}).

\subsection*{Arithmetic-to-geometric transfer}
The transfer from the arithmetic setting (automorphic form over $\Q$) to the
geometric setting (sheaf on moduli of local systems) is itself the central
open problem.  Scholze's Bourbaki lecture 1252 (March 2026)~\cite{Scholze2026}
proposes Conjecture~1.5, a ``template for a conjecture'' extending GLC to
number fields via ``Gestalts of Langlands parameters,'' with the explicit
warning: \emph{``genuinely new ideas are still needed in the global number
field case.''}  The arithmetic-to-geometric interface is further surveyed
in Zhu's 2025 ICM preprint~\cite{Zhu2025}.

The four-layer M55 uniqueness analysis (conductor minimality, Steinberg edge
$\varepsilon_{2} = -1$, twist obstruction, $\Omega$-independent ladder
rationality) precisely characterizes $\mathbf{4.5.b.a}$ within the class of
CM weight-$5$ forms over~$\Q$.  The rationality $R(f)=6/5$ is a Deligne
critical-period fact (Theorem~\ref{thm:hurwitz-bernoulli}); its categorical
interpretation in $K_0$ is conditional on Scholze~\cite{Scholze1252}
Conjecture~1.5.  The Conjecture~\ref{conj:R3C1} translation
table is:

\begin{center}
\begin{tabular}{@{}ll@{}}
\toprule
Classical (automorphic) & Spectral (Gestalt, Scholze Conj.~1.5) \\
\midrule
$K = \Q(i)$, discriminant $-4$ & dihedral image, restricted-from-$\Q(i)$ \\
Conductor $(1+i)^{2}$ & Artin conductor at $p=2$, minimal \\
Steinberg $\varepsilon_{2} = -1$ & Frobenius semisimple at $p=2$, eigenvalue $-4$ \\
$R(f) = 6/5 \in \Q$ (Deligne critical) & $K_{0}$ identity \eqref{eq:K0identity} (conditional) \\
\bottomrule
\end{tabular}
\end{center}

This translation is a restatement, not a new theorem.

%% -------------------------------------------------------------------
\section{Falsifier protocol}
\label{sec:falsifier}
%% -------------------------------------------------------------------

\subsection*{Stage 1 (completed, 2026-05-06)}

\paragraph{Stage 1 high-precision cross-check.}
The numerical identity $\pi \cdot L(f,1)/L(f,2) = 6/5$ has been verified to 80-digit precision in PARI/GP (M52 F2 v6) and independently to 50-digit precision in SageMath~10.7 (Stage 1). On 2026-05-06, an independent PARI/GP \texttt{lfunmf} subprocess bridge (M82 Method B) was executed at 75-digit precision on a separate computational pipeline, confirming
\[
\pi \cdot L(f,1)/L(f,2) = 1.20000\ldots,\qquad |R(f) - 6/5| < 1.4 \times 10^{-77},
\]
i.e.\ the identity is numerically tight to better than 77 decimal digits. This is well beyond the 50-digit Stage 1 target and constitutes essentially exact verification at any practical precision.

\textit{SageMath cross-check at $1.0\times 10^{-16}$.}  Using
\verb|Newforms(Gamma1(4),5)| to identify $f$ via $a_{2} = -4$ and the
built-in \verb|f.lseries()| engine, one obtains $|R(f) - 6/5| < 1.1\times
10^{-16}$ at Sage's default 53-bit precision.  This constitutes an independent
verification of the PARI result.  \textbf{Status: PASS (M82 Method B: $1.4\times 10^{-77}$, 2026-05-06).}

\subsection*{Stage 2 (open, estimated 50--100 CPU-hours)}
\textit{Eisenstein-symbol regulator verification.}
Using the Kato--Loeffler--Zerbes (KLZ) Eisenstein-symbol machinery
\cite{KLZ2017}, construct explicit classes
$\xi_{j}^{\mathrm{Eis}} \in H^{2}_{\mathcal{M}}(M_{f},\Q(j+1))$ for $j=1,2$
in the motivic cohomology of the motive $M_{f}$, via modular units.  Compute
the Beilinson--Deligne pairings
\[
   \mathrm{reg}_{j} \;:=\; \langle \xi_{j}^{\mathrm{Eis}}, \omega_{f}\rangle_{\mathrm{BDS}}
   \;/\; \Omega_{f}
\]
numerically (see Brunault~\cite{Brunault} for relevant explicit regulator
computations) and verify:
\[
   \frac{\mathrm{reg}_{1}}{\mathrm{reg}_{2}} \;\stackrel{?}{=}\; \frac{6}{5}
\]
to 30-digit precision.

\textbf{Explicit falsifier}: any deviation $|\mathrm{reg}_{1}/\mathrm{reg}_{2} - 6/5| > 10^{-20}$
at 30-digit verified precision kills Conjecture~R3-C-1.

\textbf{Collaborators required}: Loeffler--Zerbes (for KLZ code on CM forms),
Brunault (for Beilinson--Deligne explicit pairing).  Outreach recommended before
Stage 2 execution.

%% -------------------------------------------------------------------
\section{Bibliography}
\label{sec:bib}
%% -------------------------------------------------------------------

\begin{thebibliography}{99}

\bibitem{Beilinson1985}
A.~A.~Beilinson,
\emph{Higher regulators and values of $L$-functions},
J.\ Soviet Math.\ \textbf{30} (1985), 2036--2070.
DOI: 10.1007/BF02105861

\bibitem{Del79}
P.~Deligne,
\emph{Valeurs de fonctions $L$ et p\'eriodes d'int\'egrales},
Proc.\ Sympos.\ Pure Math.\ \textbf{33}, part~2,
Amer.\ Math.\ Soc.\ (1979), 313--346.

\bibitem{Brunault}
F.~Brunault,
\emph{Beilinson--Kato elements and modular regulators},
[Several papers, 2004--2020; specific reference to be confirmed against
CrossRef before submission.  Named here as a primary developer of explicit
Beilinson-regulator methods for modular curves.]

\bibitem{ChowlaSelberg1967}
S.~Chowla and A.~Selberg,
\emph{On Epstein's zeta-function},
J.\ Reine Angew.\ Math.\ \textbf{227} (1967), 86--110.
DOI: 10.1515/crll.1967.227.86

\bibitem{Damerell1970}
R.~M.~Damerell,
\emph{$L$-functions of elliptic curves with complex multiplication, I},
Acta Arith.\ \textbf{17} (1970), 287--301.
DOI: 10.4064/aa-17-3-287-301

\bibitem{Damerell1971}
R.~M.~Damerell,
\emph{$L$-functions of elliptic curves with complex multiplication, II},
Acta Arith.\ \textbf{19} (1971), 311--317.
DOI: 10.4064/aa-19-3-311-317

\bibitem{GLC1}
D.~Gaitsgory and S.~Raskin,
\emph{Proof of the geometric Langlands conjecture~I},
arXiv:2405.03599 (2024).

\bibitem{GLC2}
D.~Arinkin, D.~Beraldo, J.~Campbell, L.~Chen, J.~Faergeman,
D.~Gaitsgory, K.~Lin, S.~Raskin and N.~Rozenblyum,
\emph{Proof of the geometric Langlands conjecture~II},
arXiv:2405.03648 (2024).

\bibitem{GLC3}
D.~Arinkin, D.~Beraldo, J.~Campbell, L.~Chen, J.~Faergeman,
D.~Gaitsgory, K.~Lin, S.~Raskin and N.~Rozenblyum,
\emph{Proof of the geometric Langlands conjecture~III},
arXiv:2409.07051 (2024).

\bibitem{GLC4}
D.~Arinkin, D.~Beraldo, J.~Campbell, L.~Chen, J.~Faergeman,
D.~Gaitsgory, K.~Lin, S.~Raskin and N.~Rozenblyum,
\emph{Proof of the geometric Langlands conjecture~IV},
arXiv:2409.08670 (2024).

\bibitem{GLC5}
D.~Gaitsgory and S.~Raskin,
\emph{Proof of the geometric Langlands conjecture~V:
  multiplicity one},
arXiv:2409.09856 (2024).

\bibitem{Katz1976}
N.~M.~Katz,
\emph{$p$-adic interpolation of real analytic Eisenstein series},
Ann.\ Math.\ \textbf{104} (1976), 459--571.
DOI: 10.2307/1970966

\bibitem{KLZ2017}
G.~Kings, D.~Loeffler and S.~L.~Zerbes,
\emph{Rankin--Eisenstein classes and explicit reciprocity laws},
Camb.\ J.\ Math.\ \textbf{5} (2017), 1--122.
DOI: 10.4310/CJM.2017.v5.n1.a1

\bibitem{LMFDB}
The LMFDB Collaboration,
\emph{The $L$-functions and Modular Forms Database},
\url{https://www.lmfdb.org} (2013--2026).

\bibitem{SchappacherScholl1988}
N.~Schappacher and A.~J.~Scholl,
\emph{Beilinson's theorem on modular curves},
in \emph{Beilinson's Conjectures on Special Values of $L$-Functions},
Perspect.\ Math.\ \textbf{4}, Academic Press (1988), 273--304.

\bibitem{Scholl1998}
A.~J.~Scholl,
\emph{An introduction to Kato's Euler systems},
in \emph{Galois Representations in Arithmetic Algebraic Geometry},
London Math.\ Soc.\ Lecture Note Ser.\ \textbf{254}, 379--460,
Cambridge Univ.\ Press, 1998.

\bibitem{Scholze1252}
P.~Scholze,
\emph{The geometric local Langlands conjecture for groups, motives, and
condensed mathematics},
S\'em.\ Bourbaki Exp.~1252 (anticipated 2026).
[Number-field template: Conjecture~1.5.  Explicit caveat: ``genuinely new
ideas are still needed in the global number field case.'']

\bibitem{Scholze2026}
P.~Scholze,
\emph{The geometric Langlands conjecture (after Gaitsgory et al.)},
S\'eminaire Bourbaki, expos\'e 1252, March 2026.
[See also \cite{Scholze1252} for the number-field template.]

\bibitem{Shimura1976}
G.~Shimura,
\emph{The special values of the zeta functions associated with cusp forms},
Comm.\ Pure Appl.\ Math.\ \textbf{29} (1976), 783--804.
DOI: 10.1002/cpa.3160290618

\bibitem{Zhu2025}
X.~Zhu,
\emph{Arithmetic and geometric Langlands program},
arXiv:2504.07502 (2025).
[ICM 2026 survey, arithmetic-to-geometric transfer.]

\end{thebibliography}

\end{document}
